\documentclass[11pt,a4paper]{article}

% --- Paket ---
\usepackage[utf8]{inputenc}
\usepackage[T1]{fontenc}
\usepackage[swedish]{babel}
\usepackage{amsmath, amssymb, amsthm}
\usepackage{mathtools}
\usepackage{geometry}
\geometry{margin=2.5cm}
\usepackage{parskip}
\usepackage{xcolor}
\usepackage{hyperref}
\hypersetup{
    colorlinks=true,
    linkcolor=blue!50!black,
    urlcolor=blue!50!black
}

% --- Egna kommandon ---
\newcommand{\R}{\mathbb{R}}
\newcommand{\E}{\mathbb{E}}
\newcommand{\Var}{\operatorname{Var}}
\newcommand{\dd}{\mathrm{d}}
\DeclareMathOperator{\Cov}{Cov}

% --- Satsmiljöer ---
\theoremstyle{plain}
\newtheorem{theorem}{Sats}
\newtheorem{lemma}[theorem]{Lemma}

\theoremstyle{definition}
\newtheorem{definition}[theorem]{Definition}

\theoremstyle{remark}
\newtheorem*{remark}{Anmärkning}

% --- Titel ---
\title{Härledning av obligationspriset i Vasicek-modellen \\ via affin ansats}
\author{Projektarbete SF1693}
\date{\today}

\begin{document}

\maketitle

\section{Problemformulering}

I Vasicek-modellen \cite{Vasicek1977} antas korträntan $\{r_t\}_{t\geq 0}$ under det riskneutrala måttet $\mathbb{Q}$ följa en Ornstein--Uhlenbeck-process,
\begin{equation}
    \dd r_t = \theta(\mu - r_t)\,\dd t + \sigma\,\dd W_t, \qquad r_0 \in \R,
    \label{eq:vasicek-sde}
\end{equation}
där $\theta>0$ är återgångshastigheten, $\mu\in\R$ den långsiktiga medelnivån, $\sigma>0$ volatiliteten, och $\{W_t\}_{t\geq0}$ en $\mathbb{Q}$-Wienerprocess. Priset vid tid $t$ på en nollkupongsobligation med förfall $T\geq t$ ges av den riskneutrala värderingsformeln
\begin{equation}
    P(t,T) \;=\; \E^{\mathbb{Q}}\!\left[\, \exp\!\left(-\!\int_t^T r_s\,\dd s\right) \,\Big|\, r_t = r \right].
    \label{eq:bond-def}
\end{equation}
Målet är att härleda en sluten formel för $P(t,T)$ genom att utnyttja Vasicek-modellens \emph{affina} struktur.

\section{Förberedande räkneregler}

\begin{lemma}[Feynman--Kac]
\label{lem:fk}
Antag att $r_t$ uppfyller \eqref{eq:vasicek-sde} och att $F\in C^{2,1}(\R\times[0,T])$ uppfyller den parabolska PDE:n
\[
    \frac{\partial F}{\partial t} + \theta(\mu - r)\frac{\partial F}{\partial r} + \tfrac12\sigma^2\frac{\partial^2 F}{\partial r^2} - r\,F \;=\; 0,
\]
med terminalvillkor $F(T,r) = 1$ och tillräcklig integrabilitet. Då gäller
\[
    F(t,r) \;=\; \E^{\mathbb{Q}}\!\left[\, e^{-\int_t^T r_s\,\dd s} \,\Big|\, r_t = r \right] \;=\; P(t,T).
\]
\end{lemma}

\begin{remark}
Resultatet följer av Itôs formel applicerad på processen $M_s = e^{-\int_t^s r_u\,\dd u}\,F(s,r_s)$: PDE:n säger precis att drifttermen i $\dd M_s$ försvinner, så $M$ är en martingal. Identifikationen $\E^{\mathbb{Q}}[M_T \mid \mathcal{F}_t] = M_t$ tillsammans med $F(T,r)=1$ ger formeln.
\end{remark}

Vi kommer alltså att lösa PDE:n istället för att beräkna den betingade förväntan i \eqref{eq:bond-def} direkt.

\section{Affin ansats}

\subsection*{Steg 1: Ansats}

Eftersom drifttermen $\theta(\mu-r)$ är linjär i $r$ och diffusionen $\sigma^2$ är konstant tillhör Vasicek-modellen klassen \emph{affina termstrukturmodeller}. För sådana modeller är logaritmen av obligationspriset linjär i $r$. Vi gör därför ansatsen
\begin{equation}
    P(t,T,r) \;=\; \exp\!\bigl(A(t,T) - B(t,T)\,r\bigr),
    \label{eq:affine-ansatz}
\end{equation}
med terminalvillkoren $A(T,T) = 0$ och $B(T,T) = 0$, så att $P(T,T,r) = 1$ uppfylls för alla $r$.

\subsection*{Steg 2: Insätt ansatsen i PDE:n}

De partiella derivatorna av \eqref{eq:affine-ansatz} är
\[
    \frac{\partial P}{\partial t} = \bigl(A_t - B_t\,r\bigr)P, \qquad
    \frac{\partial P}{\partial r} = -B\,P, \qquad
    \frac{\partial^2 P}{\partial r^2} = B^2\,P,
\]
där $A_t \coloneqq \partial A/\partial t$ och analogt för $B_t$. Insättning i Feynman--Kac-PDE:n och division med den positiva faktorn $P$ ger
\begin{equation}
    A_t - B_t\,r \;-\; \theta(\mu-r)\,B \;+\; \tfrac12\sigma^2 B^2 \;-\; r \;=\; 0.
    \label{eq:after-substitution}
\end{equation}

\subsection*{Steg 3: Separation efter potenser av $r$}

Samla termer i ekvation \eqref{eq:after-substitution} efter potenser av $r$:
\begin{equation}
    \underbrace{\bigl[\,-B_t + \theta B - 1\,\bigr]}_{\text{koefficient för }r}\, r \;+\; \underbrace{\bigl[\,A_t - \theta\mu\,B + \tfrac12\sigma^2 B^2\,\bigr]}_{\text{konstant i }r} \;=\; 0.
    \label{eq:r-separation}
\end{equation}
Eftersom \eqref{eq:r-separation} skall hålla för \emph{alla} $r$ måste varje hakparentes vara identiskt noll. Detta ger ett par frikopplade ordinära differentialekvationer i $t$ med terminalvillkor:
\begin{align}
    B_t(t,T) &= \theta\,B(t,T) - 1, & B(T,T) &= 0, \label{eq:B-ode}\\
    A_t(t,T) &= \theta\mu\,B(t,T) - \tfrac12\sigma^2 B(t,T)^2, & A(T,T) &= 0. \label{eq:A-ode}
\end{align}
Notera att \eqref{eq:B-ode} är fristående och kan lösas först; därefter ger \eqref{eq:A-ode} $A$ via direkt integration.

\section{Lösning av ODE-systemet}

\subsection*{Lösning av $B$}

Inför löptiden $\tau \coloneqq T - t$ och betrakta $B$ som funktion av $\tau$ (under fixerat $T$). Då är $\dd/\dd\tau = -\dd/\dd t$, så \eqref{eq:B-ode} blir
\begin{equation}
    \frac{\dd B}{\dd\tau} \;=\; 1 - \theta B, \qquad B(0) = 0.
    \label{eq:B-tau}
\end{equation}
Detta är en linjär första ordningens ODE. Den integrerande faktorn $e^{\theta\tau}$ ger
\[
    \frac{\dd}{\dd\tau}\bigl(e^{\theta\tau}B\bigr) = e^{\theta\tau},
\]
så $e^{\theta\tau}B(\tau) = (e^{\theta\tau} - 1)/\theta$, dvs.
\begin{equation}
    \boxed{\,B(\tau) \;=\; \dfrac{1 - e^{-\theta\tau}}{\theta}.\,}
    \label{eq:B-solution}
\end{equation}

\subsection*{Lösning av $A$}

Med $B$ känd integreras \eqref{eq:A-ode} från $t$ till $T$. Eftersom $A(T,T)=0$:
\begin{equation}
    A(t,T) \;=\; -\!\int_t^T A_s(s,T)\,\dd s \;=\; \int_0^\tau \!\Bigl[\tfrac12\sigma^2 B(u)^2 - \theta\mu\,B(u)\Bigr]\,\dd u,
    \label{eq:A-integral}
\end{equation}
där substitutionen $u = T - s$ använts. Vi behöver alltså två integraler.

\paragraph{Första integralen.}
\begin{align}
    \int_0^\tau B(u)\,\dd u
    &= \frac{1}{\theta}\int_0^\tau \bigl(1 - e^{-\theta u}\bigr)\,\dd u
    = \frac{1}{\theta}\!\left[\,u + \frac{e^{-\theta u}}{\theta}\,\right]_0^\tau \nonumber\\
    &= \frac{1}{\theta}\!\left[\,\tau - \frac{1 - e^{-\theta\tau}}{\theta}\,\right]
    = \frac{\tau - B(\tau)}{\theta}.
    \label{eq:int-B}
\end{align}

\paragraph{Andra integralen.} Med $B(u)^2 = \theta^{-2}\bigl(1 - 2e^{-\theta u} + e^{-2\theta u}\bigr)$ är
\begin{align}
    \int_0^\tau B(u)^2\,\dd u
    &= \frac{1}{\theta^2}\!\left[\,\tau - 2B(\tau) + \frac{1 - e^{-2\theta\tau}}{2\theta}\,\right].
    \label{eq:int-B2-raw}
\end{align}
Genom att skriva $1 - e^{-2\theta\tau} = (1 - e^{-\theta\tau})(1 + e^{-\theta\tau}) = \theta B(\tau)\bigl(2 - \theta B(\tau)\bigr)$ förenklas sista termen till
\[
    \frac{1 - e^{-2\theta\tau}}{2\theta} \;=\; B(\tau) - \frac{\theta\,B(\tau)^2}{2},
\]
så att
\begin{equation}
    \int_0^\tau B(u)^2\,\dd u \;=\; \frac{1}{\theta^2}\!\left[\,\tau - B(\tau) - \frac{\theta\,B(\tau)^2}{2}\,\right].
    \label{eq:int-B2}
\end{equation}

\paragraph{Sätt ihop.} Insättning av \eqref{eq:int-B} och \eqref{eq:int-B2} i \eqref{eq:A-integral} ger
\begin{align}
    A(t,T)
    &= \frac{\sigma^2}{2\theta^2}\!\left[\,\tau - B(\tau) - \frac{\theta\,B(\tau)^2}{2}\,\right] - \mu\bigl(\tau - B(\tau)\bigr) \nonumber\\[2pt]
    &= \bigl(B(\tau) - \tau\bigr)\!\left[\,\mu - \frac{\sigma^2}{2\theta^2}\,\right] - \frac{\sigma^2 B(\tau)^2}{4\theta}.
\end{align}
Med beteckningen för den \emph{långsiktiga räntan}
\begin{equation}
    R_\infty \;\coloneqq\; \mu - \frac{\sigma^2}{2\theta^2},
    \label{eq:Rinf}
\end{equation}
får vi den kompakta formeln
\begin{equation}
    \boxed{\,A(t,T) \;=\; \bigl(B(\tau) - \tau\bigr)\,R_\infty \;-\; \frac{\sigma^2 B(\tau)^2}{4\theta}.\,}
    \label{eq:A-solution}
\end{equation}

\section{Resultat}

\begin{theorem}[Vasicek-formeln för obligationspriset]
\label{thm:vasicek}
Med $r_t$ definierad genom \eqref{eq:vasicek-sde} och $\tau = T-t$ gäller
\begin{equation}
    P(t,T,r) \;=\; \exp\!\Bigl(\,A(t,T) - B(t,T)\,r\,\Bigr),
    \label{eq:vasicek-price}
\end{equation}
där
\[
    B(t,T) = \frac{1 - e^{-\theta\tau}}{\theta}, \qquad
    A(t,T) = \bigl(B(t,T) - \tau\bigr) R_\infty - \frac{\sigma^2 B(t,T)^2}{4\theta},
\]
och $R_\infty = \mu - \sigma^2/(2\theta^2)$.
\end{theorem}

\begin{remark}
Att $\log P$ är linjär i $r$ är just innebörden i att Vasicek är en \emph{affin termstrukturmodell}. Detta var inte en lycklig gissning utan en konsekvens av den affina drift--diffusionsstrukturen i \eqref{eq:vasicek-sde}: ansatsen \eqref{eq:affine-ansatz} reducerar PDE:n till två ODE:n i koefficienterna, vilket är dramatiskt enklare än den ursprungliga PDE:n.
\end{remark}

\section{Räntestruktur och gränsvärden}

Den kontinuerligt sammansatta avkastningen mellan $t$ och $T$ definieras som $Y(t,T) = -\log P(t,T)/\tau$. Av \eqref{eq:vasicek-price} följer
\begin{equation}
    Y(t,T) \;=\; \frac{B(t,T)}{\tau}\,r \;-\; \frac{A(t,T)}{\tau},
    \label{eq:yield}
\end{equation}
en \emph{affin} funktion av kortränta $r$. Två kontroller:
\begin{itemize}
    \item \textbf{Kort horisont:} När $\tau\to 0^+$ ger Taylor-utveckling $B(\tau) = \tau - \theta\tau^2/2 + O(\tau^3)$ och $A(t,T) = O(\tau^3)$, så $Y(t,T) \to r$. Avkastningen i den närmaste framtiden sammanfaller alltså med korträntan, som väntat.
    \item \textbf{Lång horisont:} När $\tau\to\infty$ är $B(\tau)\to 1/\theta$ medan $-A(t,T)/\tau \to R_\infty$, så $Y(t,T) \to R_\infty$ oberoende av $r$. Detta motiverar namnet ``långsiktig ränta'' för $R_\infty$ i \eqref{eq:Rinf}.
\end{itemize}

\section{Sammanfattning}

Genom den affina ansatsen $P = e^{A - Br}$ reducerades den parabolska PDE:n från Feynman--Kac till två frikopplade ODE:n i koefficienterna $A$ och $B$:
\[
    B_t = \theta B - 1, \qquad A_t = \theta\mu\,B - \tfrac12\sigma^2 B^2,
\]
båda med terminalvärde noll vid $t=T$. Lösningarna är
\[
    B(\tau) = \frac{1 - e^{-\theta\tau}}{\theta}, \qquad
    A(t,T) = \bigl(B(\tau) - \tau\bigr)R_\infty - \frac{\sigma^2 B(\tau)^2}{4\theta},
\]
med $R_\infty = \mu - \sigma^2/(2\theta^2)$. Vasiceks slutna formel
\[
    P(t,T,r) = \exp\!\bigl(A(t,T) - B(t,T)\,r\bigr)
\]
ger en analytisk referenslösning som det numeriska schemat för motsvarande PDE kan valideras mot.

\begin{thebibliography}{9}
\bibitem{Vasicek1977}
O.~Vasicek, \emph{An equilibrium characterization of the term structure}, Journal of Financial Economics \textbf{5} (1977), 177--188.
\end{thebibliography}

\end{document}
